%! TEX root = ../main.tex
\documentclass[main]{subfiles}

\begin{document}
\chapter{今後の展望}
本研究では、DCASE2024 Task 4のベースラインモデルを基盤とし、半教師あり学習手法Mean Teacher（MT）の改良、ドメイン汎化、及び後処理の改善を統合することで、音響イベント検出の精度向上を図った。 
結果として、負例を活用したCMT（$CMT_{scale}$）は、従来のMTや単純なCMTと比較して精度向上を達成した。
また、周波数ごとの重要度に基づくMixStyleの統合が有効であることを再確認し、軽量なCNNを用いた場合でもドメイン汎化の効果が得られることを示した。\par

今後の課題として、CMTとMixStyleを併用した際に観測された、ベストモデル到達エポック数の減少（早期収束または不安定性）に関する詳細な解析や、
他の半教師あり学習手法との性能比較が挙げられる。
また、実用化の観点からはモデルの軽量化が急務である。MTは教師モデルを保持する必要があるため、パラメータ数が増大する傾向にある。
本研究では既存の標準的手法であるMTの改良に焦点を当てたが、エッジデバイスでの見守りシステムへの実装を考慮すると、よりパラメータ効率の良い手法の検討が必要である。\par 

さらに、音響イベント検出は「いつ」「何が」鳴ったかを同時に推定するマルチタスク性を持つため、モデルが肥大化しやすい。
エッジデバイスでの動作を前提とするならば、機械学習モデルの役割をクラス分類に限定し、時間的な区間推定は統計的な手法に委任するなど、
タスクとモデル構造の簡素化を図るアプローチも有効な解決策になると考えられる。

\end{document}